\section{Stars and Bars}

\textbf{Problema base.} Quantas soluções inteiras tem
\[
x_1 + x_2 + \dots + x_k = n, \qquad x_i \ge 0 \;?
\]

\textbf{Ideia.} Represente $n$ como $n$ estrelas ($\star$) e separe-as em $k$ grupos
usando $k-1$ barras ($|$). Cada arrumação de $n$ estrelas e $k-1$ barras em fila
corresponde a exatamente uma solução (o grupo $i$ é a quantidade de estrelas
entre a barra $i-1$ e a barra $i$). Ex.: $n=5$, $k=3$: \;
$\star\star|\star|\star\star \;\leftrightarrow\; (2,1,2)$ \; e \;
$||\star\star\star\star\star \;\leftrightarrow\; (0,0,5)$.

Basta escolher as posições das barras entre $n+k-1$ posições:
\[
\boxed{\binom{n+k-1}{k-1}}
\]

\subsection*{Variações}

\begin{itemize}
  \item \textbf{Variáveis positivas} ($x_i \ge 1$): coloque as barras nos $n-1$ espaços
        entre estrelas, no máximo uma por espaço:
        \[ \binom{n-1}{k-1}. \]

  \item \textbf{Limite inferior} ($x_i \ge l_i$): substitua $y_i = x_i - l_i \ge 0$.
        Recai no caso base com $n' = n - \sum l_i$ (se $n' < 0$, a resposta é $0$):
        \[ \binom{n' + k - 1}{k-1}. \]

  \item \textbf{Soma com desigualdade} ($x_1 + \dots + x_k \le n$, $x_i \ge 0$):
        adicione uma variável de folga $x_{k+1} = n - \sum x_i \ge 0$. Agora são $k+1$ variáveis
        com soma exatamente $n$:
        \[ \binom{n+k}{k}. \]

  \item \textbf{Limite superior} ($0 \le x_i \le u_i$): inclusão-exclusão. Para um conjunto $S$
        de variáveis que \emph{violam} o limite, force $x_i \ge u_i + 1$ para $i \in S$:
        \[
        \sum_{S \subseteq \{1,\dots,k\}} (-1)^{|S|}
        \binom{n - \sum_{i \in S}(u_i+1) + k - 1}{k-1},
        \]
        considerando o termo $0$ quando o topo fica negativo.

  \item \textbf{Mesmo limite superior} ($0 \le x_i \le c$ para todo $i$): os conjuntos $S$ de
        mesmo tamanho contribuem igual, então
        \[
        \sum_{j=0}^{k} (-1)^j \binom{k}{j} \binom{n - j(c+1) + k - 1}{k-1},
        \]
        parando quando $n - j(c+1) < 0$.
\end{itemize}

\subsection*{Problemas equivalentes}

A mesma fórmula $\binom{n+k-1}{k-1} = \binom{n+k-1}{n}$ conta:
\begin{itemize}
  \item distribuir $n$ objetos \textbf{idênticos} em $k$ caixas \textbf{distintas};
  \item \textbf{multiconjuntos} de tamanho $n$ com elementos de $k$ tipos
        (combinação com repetição);
  \item sequências \textbf{não-decrescentes} de tamanho $n$ com valores em $[1, k]$
        ($x_i$ = quantas vezes o valor $i$ aparece);
  \item monômios de grau $n$ em $k$ variáveis.
\end{itemize}
